\scanpage{68}
\deleted{Ceci posé} Cette relation est évidemment \deleted{symétrique}\inserted{réflexive}, et
elle est transitive grâce au fait que
\[
\pi_0(\Sigma_{(2)})\longrightarrow
(\pi_0(\Sigma_{(1)}),\pi_0(\uline{b}))\times_{\pi_0(\Sigma_{(0)})}
(\pi_0(\Sigma_{(1)}),\pi_0(\uline{s}))
\]
est surjective, ce qui est facile à prouver,
grâce à l'hypothèse que $\uline{b}$ est fibrant, par le
\uline{Lemme} : Soit
\diagram{\begin{tikzcd}
X\arrow[d,"f"'] & X'\arrow[l,"f'"']\arrow[d,"g'"]\\
Y & Y'\arrow[l,"g"]
\end{tikzcd}}
un carré
cartésien d'espaces top., avec $f$ fibrant,
alors $\pi_0(X')\longrightarrow\pi_0(X)\times_{\pi_0(Y)}\pi_0(Y')$ est surjectif.
Ceci posé, considérons la catégorie top. discrète associée
à l'ens \inserted{pré}ordonné $I_0$\inserted{$=\pi_0(\Sigma_0)$} que par abus de langage
\inserted{[\unsure{morphisme|anglicisme})}
on désignera par $I$. On a alors un foncteur
\inserted{canonique} de catégories topologiques \deleted{canonique}
\[
\uline{\Sigma}\xrightarrow{\uline{\varphi}}I
\]
qui est donc donné par des applications continues
\[
\left\{\begin{aligned}
\Sigma_{(0)}&\xrightarrow{\varphi_0}I_0=I\\
\Sigma_{(1)}&\xrightarrow{\varphi_1}I_1=\text{graphe de la relation d'ordre}
\end{aligned}\right.
\]
\inserted{$=\text{image de }\pi_0(\Sigma_1)$ \deleted{dans $\unsure{\pi_0(\Sigma_1)\times I}$}
dans $\pi_0(\Sigma_0)\times\pi_0(\Sigma_0)=I$}
i.e. des décompositions en sommes directes
\[
\Sigma_{(0)}=\coprod_{i\in I}\Sigma_i
\]
\[
\Sigma_{(1)}=\deleted{\sum}\coprod_{\substack{i,j\in I\\i\leq j}}\Sigma_{ij}
\]
\margin{\uline{NB} on suppose les comp. connexes des $\Sigma_0$ et $\Sigma_1$ ouvertes}
compatibles avec $\uline{s},\uline{b}$
i.e. $\uline{s}:\Sigma_{ij}\longrightarrow\Sigma_i$
$\uline{b}:\Sigma_{ij}\longrightarrow\Sigma_j$
avec $\Sigma_{(0)}\xrightarrow{\delta}\Sigma_{(1)}$ ($\delta(x)=\deleted{\unsure{\Sigma}}\operatorname{id}^{\uline{\Sigma}}_x$)
appliquant $\Sigma_i$ dans $\Sigma_{ii}$, \deleted{\unsure{et ainsi}} ce qui
implique déjà la compatibilité de $\varphi=(\varphi_0,\varphi_1)$ avec
la composition). De plus, $\varphi$ est un
\mspage{67}
